% ================================================================
\partdivider{A}{Mechanical Properties of Matter}{How solids resist being deformed and how fluids push, float and flow: the mechanics of bulk matter, built from the forces and energy of Volume~I.}{\chref{ch:solids} Mechanical Properties of Solids\\[1.2mm]
\chref{ch:fluids} Mechanical Properties of Fluids}
\ncertchapter{1}{Mechanical Properties of Solids}\chlabel{ch:solids}
% ================================================================

Why can a steel cable hold up a bridge while a rubber band of the same thickness cannot? Why does a
diving board bend under you and spring back, while a biscuit snaps? Why are bones hollow, girders
I-shaped, and aircraft made of aluminium alloy rather than lead? Each of these is a question about
\textbf{elasticity}: how a solid deforms when forces act on it, and whether it returns to its shape when
they are removed.

Volume~I mostly avoided the question. Its rigid bodies (\chref{ch:rotation}) do not deform at all, and
its springs (\chref{ch:connected}) obey Hooke's law, $F=-kx$, with a spring constant $k$ that belongs to
one particular spring. Neither idea is enough to design a bridge. A real cable stretches, and an engineer
needs to predict how far before the cable is made, from the properties of steel rather than of one
finished cable. This chapter supplies that prediction. It rests on two ideas, \emph{stress} and
\emph{strain}, which remove the size and shape of the sample, tied together by one number for each
material and each way of deforming it, its \emph{modulus}. Hooke's spring constant turns out to be a
modulus multiplied by a ratio of dimensions.

The limits deserve stating from the start. Everything here assumes small deformations, a material that
is the same in every direction, and loads applied slowly enough that nothing vibrates. Beyond the
elastic limit the simple laws fail, and the chapter describes, without a full theory, what happens
there: permanent deformation and, eventually, fracture.

\begin{motivbox}{Why elasticity matters}
Every structure trusted with a life is an elasticity calculation. A lift cable, a bridge span and a hip
implant must neither stretch too far nor break, which means keeping the stress (force per area) safely
below the limit of the material: that is why cables are thick and why the same load is safe on one
material and dangerous on another. Shape matters as much as material: a hollow bone, an I-shaped girder
and a tube resist bending far better than a solid rod of the same mass, because bending is resisted most
by material far from the axis, an idea that echoes the moment of inertia of \chref{ch:rotation}. And
elastic energy, stored and released, is what lets a diving board spring, a car suspension absorb a bump
and a pole-vaulter's pole hurl an athlete over the bar; released suddenly, the same energy makes a
snapped cable dangerous.
\end{motivbox}

\section{Stress and Strain}\label{sec:solids-stress}

When a rod is pulled, two things compete. The external force tries to stretch it; the bonds between its
atoms pull back, trying to restore their natural spacing. In equilibrium these balance, and the rod
settles at a slightly greater length. To describe this in a way that does not depend on the particular
rod, we need quantities that do not change when the rod is made thicker or longer: force spread over
area, and stretch relative to length.

\begin{defbox}{Stress and strain}
Take a rod of natural length $\ell$ and cross-sectional area $A$, pulled at its two ends by equal and
opposite forces of magnitude $F$, so that it stretches by $\Delta\ell$.
\begin{itemize}
\item \textbf{Stress} is the internal force per unit area across any cross-section,
      \[ \sigma=\frac{F}{A}\qquad(\text{unit }\si{N/m^2}=\si{Pa}). \]
      Imagine cutting the rod anywhere along its length: the two pieces must pull on each other with
      force $F$, or each piece would accelerate. That internal pull, spread over the area of the cut, is
      the stress. It is a scalar here because the force is perpendicular to the cut; it is taken as
      positive in tension and negative in compression.
\item \textbf{Strain} is the \emph{fractional} deformation,
      \[ \epsilon=\frac{\Delta\ell}{\ell}\qquad(\text{dimensionless}). \]
\end{itemize}
Stress measures how hard the material is being worked; strain measures how much it gives.
\end{defbox}

\inwords{Dividing by area and by length describes the \emph{material}, not the particular sample. A
thick rope holds more than a thin one made of the same fibre, but at the same \emph{stress} both are
equally close to breaking. That is why engineers quote stress, not force: it is the fair measure of how
hard a material is being pushed, whatever its size. The unit of stress is the unit of pressure, and the
pressure of a fluid (\chref{ch:fluids}) is a special case of stress, one that pushes equally in every
direction.}

Pulling along the length is only one way to deform a solid. The direction of the force relative to the
surface it acts on decides what kind of deformation follows, and there are three kinds.

\begin{defbox}{Three kinds of stress and strain}
The three are shown in Figure~\ref{fig:solids-stresskinds}.
\begin{itemize}
\item \textbf{Tensile or compressive.} The force is \emph{perpendicular} to the face it acts on,
      pulling out or pushing in along the length. The response is a \textbf{longitudinal strain}
      $\Delta\ell/\ell$. Examples: a stretched wire, a loaded pillar.
\item \textbf{Shear.} The force is \emph{parallel} to the face (tangential), and an equal and opposite
      tangential force acts on the opposite face, so the top of a block slides over its bottom without
      the block as a whole moving. The shearing stress is the tangential force per area of the face,
      $F/A$. The response is a \textbf{shear strain}, the angle $\theta$ (in radians) through which an
      originally perpendicular side tilts; if the top moves by $\Delta x$ relative to the bottom, a
      height $h$ below it, then $\tan\theta=\Delta x/h$, and for the small angles of elasticity
      $\theta\approx\Delta x/h$ (the small-angle approximation of \chref{ch:toolkit}). Examples: a twisted
      shaft, a sheet cut by scissors.
\item \textbf{Volume (hydrostatic).} An equal increase in pressure $\Delta P$ acts on \emph{every} face
      at once. The response is a \textbf{volume strain} $\Delta V/V$, the fractional change in volume.
      Example: a body sunk deep in the sea.
\end{itemize}
Each kind of strain is matched by its own elastic modulus, defined in the next section.
\end{defbox}

\begin{figure}[htbp]
\centering
\tikzfig{solids_stresskinds}
\figcaption{fig:solids-stresskinds}{The three ways to deform a solid, with the deformations exaggerated
about a hundredfold (dashed outlines are the undeformed shapes). Tensile: forces normal to the end faces
lengthen the block, which also grows slightly narrower. Shear: equal and opposite tangential forces on
the top and bottom faces tilt the sides through the angle $\theta$, which \emph{is} the strain. Volume:
the same extra pressure on every face shrinks the block without changing its shape.}
\end{figure}

\begin{exbox}{Shearing a metal block}\label{ex:solids-shear}
An aluminium cube of side \SI{10}{cm} has its bottom face glued to a bench. A tangential force of
\SI{1.0e5}{N} is applied to its top face. How far does the top face move? (Shear modulus of aluminium,
defined in the next section: $G=\SI{2.5e10}{Pa}$, so that shear strain $=$ shear stress$/G$.)

\stepn{1} Shear stress: $F/A=\num{1.0e5}/(0.10)^2=\SI{1.0e7}{Pa}$.

\stepn{2} Shear strain: $\theta=\num{1.0e7}/\num{2.5e10}=\SI{4.0e-4}{rad}$, about $0.02^\circ$.

\stepn{3} Displacement of the top: $\Delta x=h\theta=(0.10)(\num{4.0e-4})=\SI{4.0e-5}{m}$.

\smallskip
A force equal to the weight of ten tonnes moves the top face by \SI{0.04}{mm}: metals are stiff, which is
why the rigid-body model of Volume~I works so well for them. \ans{\SI{0.04}{mm}}
\end{exbox}

\section{Hooke's Law and the Elastic Moduli}\label{sec:solids-moduli}

Robert Hooke found in 1660 that, for small deformations, the extension of a spring or wire is
proportional to the force: \emph{ut tensio, sic vis}, ``as the extension, so the force.''
\chref{ch:connected} used that law for a spring. In the language of the last section it becomes a
statement about any material: for small strains, \emph{stress is proportional to strain}. The constant of
proportionality is the material's \textbf{modulus of elasticity}, and it measures the stiffness of the
material rather than of one object.

\begin{keybox}{The three elastic moduli}
For small strains, stress $=$ modulus $\times$ strain, with one modulus for each kind of deformation:
\[
\underbrace{Y=\frac{F/A}{\Delta\ell/\ell}}_{\text{Young's modulus}},\qquad
\underbrace{G=\frac{F/A}{\theta}}_{\text{shear modulus}},\qquad
\underbrace{B=-\frac{\Delta P}{\Delta V/V}}_{\text{bulk modulus}}.
\]
Each is a stress divided by a dimensionless strain, so each has the unit of pressure, the pascal. A
\emph{large} modulus means a stiff material: a large stress produces only a small strain. The bulk
modulus carries a minus sign because an increase in pressure \emph{decreases} the volume, and the sign
makes $B$ positive. Its reciprocal, $1/B$, is called the \textbf{compressibility}. (The bulk modulus is
written $B$ in this volume because $K$ is reserved for thermal conductivity in \chref{ch:thermal}.)
\end{keybox}

\inwords{One relation, stress $=$ modulus $\times$ strain, wears three hats depending on how the material
is pushed. Pull a wire and $Y$ governs it; shear or twist a shaft and $G$ governs it; squeeze from all
sides and $B$ governs it. A fluid cannot be pulled or sheared into a new shape that it then holds, so it
has no $Y$ or $G$ at all, only a bulk modulus; that is why sound in a gas, a wave of compression, depends
on the bulk modulus of the gas.}

The relation between the modulus of a material and the spring constant of an object made from it is the
first thing worth deriving, because it connects this chapter to Volume~I.

\begin{derivbox}{A wire is a spring: $k=YA/\ell$ and $\Delta\ell=F\ell/AY$}
\stepn{1} Solve the definition of Young's modulus for the force:
\[
Y=\frac{F/A}{\Delta\ell/\ell}\qquad\Longrightarrow\qquad F=\frac{YA}{\ell}\,\Delta\ell .
\]
\stepn{2} This is Hooke's law, $F=kx$ in magnitude, with extension $x=\Delta\ell$ and spring constant
\[
\boxed{\;k=\frac{YA}{\ell}\;}
\]
The spring constant of \chref{ch:connected} is therefore not a property of the material alone: it is the
modulus (the material) multiplied by $A/\ell$ (the shape). A thicker wire is a stiffer spring; a longer
one is a softer spring.

\stepn{3} Solving instead for the extension,
\[
\boxed{\;\Delta\ell=\frac{F\ell}{AY}\;}
\]
The stretch grows with the load $F$ and the length $\ell$ (a longer wire has more material to stretch),
and shrinks with the area $A$ and the stiffness $Y$. Doubling the length doubles the stretch; halving the
radius quarters the area and so \emph{quadruples} the stretch. That last consequence is the most tested
in examinations.
\end{derivbox}

\begin{exbox}{How much does a steel wire stretch?}\label{ex:solids-wire}
A steel wire \SI{1.0}{m} long with cross-sectional area \SI{1.0}{mm^2} carries a \SI{10}{kg} load. How far
does it stretch, and what is its spring constant? ($Y_{\text{steel}}=\SI{2.0e11}{Pa}$,
$g=\SI{9.8}{m/s^2}$.)

\stepn{1} Force and area: $F=mg=\SI{98}{N}$; $A=\SI{1.0e-6}{m^2}$.

\stepn{2} Stress: $\sigma=F/A=\SI{9.8e7}{Pa}=\SI{98}{MPa}$.

\stepn{3} Extension: $\Delta\ell=\dfrac{F\ell}{AY}=\dfrac{(98)(1.0)}{(\num{1.0e-6})(\num{2.0e11})}
=\SI{4.9e-4}{m}\approx\SI{0.49}{mm}$.

\stepn{4} Spring constant: $k=YA/\ell=\SI{2.0e5}{N/m}$, about a hundred times stiffer than a typical
laboratory spring.

\smallskip
Half a millimetre under a load heavier than a bag of cement: steel gives very little, which is exactly why
loads are hung from it. \ans{\SI{0.49}{mm}; \SI{2.0e5}{N/m}}
\end{exbox}

\begin{trapbox}{a longer wire is not a stronger one}
A longer wire \emph{stretches more} for the same load ($\Delta\ell\propto\ell$), but it \emph{breaks at
the same load}, because breaking depends on the stress $F/A$, which does not involve the length. Length
affects how much a wire gives, not when it fails.
\end{trapbox}

Typical values, collected in Table~\ref{tab:solids-moduli}, show how widely the moduli differ, and
they hide one more lesson: for solids the three moduli are of the same order, while for liquids and gases
only $B$ exists.

\begin{tablebox}{Approximate elastic moduli (in \si{GPa}, $\SI{1}{GPa}=\SI{e9}{Pa}$)}
\centering
\begin{tabular}{@{}lccc@{}}
\toprule
Material & $Y$ & $G$ & $B$\\
\midrule
steel & $200$ & $80$ & $160$\\
copper & $120$ & $44$ & $140$\\
aluminium & $70$ & $25$ & $70$\\
glass & $65$ & $26$ & $40$\\
bone (along its length) & $16$ & -- & --\\
soft rubber & about $0.001$ to $0.01$ & -- & $2$\\
water & -- & -- & $2.2$\\
air (at atmospheric pressure, constant temperature) & -- & -- & $0.0001$\\
\bottomrule
\end{tabular}\label{tab:solids-moduli}

\smallskip
\raggedright
Values vary by tens of per cent with composition and treatment. A dash means the modulus does not exist
(fluids) or is not needed here.
\end{tablebox}

\why{The stiffness of a material is set by the strength of its interatomic bonds, so the modulus is a
window onto the microscopic world. Steel resists stretching about $10^{5}$ times more strongly than soft
rubber, because metallic bonds are far stiffer than the tangled polymer chains of rubber, which stretch by
straightening out rather than by pulling atoms apart. Rubber is nevertheless hard to \emph{compress}
($B\approx\SI{2}{GPa}$, like water): its chains straighten easily, but its atoms resist being crowded
together. Choosing a material, steel for a bridge, rubber for a tyre, aluminium for a wing, is largely
choosing a modulus.}

\begin{exbox}{How much does water compress in the deep sea?}\label{ex:solids-water}
By what fraction is the volume of water reduced at a depth of \SI{1.0}{km}, where the pressure exceeds the
surface pressure by $\rho gh$? ($B_{\text{water}}=\SI{2.2e9}{Pa}$, $\rho=\SI{1000}{kg/m^3}$.)

\stepn{1} $\Delta P=\rho gh=(1000)(9.8)(1000)=\SI{9.8e6}{Pa}$, about $97$ atmospheres.

\stepn{2} $\Delta V/V=-\Delta P/B=-\num{9.8e6}/\num{2.2e9}=-\num{4.5e-3}$.

\smallskip
Less than half a per cent, under nearly a hundred atmospheres. This is why \chref{ch:fluids} treats
liquids as incompressible, and why that model fails only in the deepest oceans or at very high pressures.
\ans{$0.45\%$ smaller}
\end{exbox}

\subsection{Poisson's ratio: stretch one way, thin the other}

A stretched rubber band visibly grows narrower as it lengthens. Every solid does the same, though usually
too little to see: pulling it longer makes it thinner, and squashing it shorter makes it bulge. The two
strains are linked by a pure number.

\begin{defbox}{Poisson's ratio}
When a rod of width (or diameter) $d$ is stretched, its \textbf{longitudinal strain} $\Delta\ell/\ell$ is
accompanied by a \textbf{lateral strain} $\Delta d/d$ of the opposite sign. Their ratio defines
\textbf{Poisson's ratio},
\[
\boxed{\;\nu=-\frac{\Delta d/d}{\Delta\ell/\ell}\;}
\]
The minus sign makes $\nu$ positive. It is dimensionless and, for most solids, between about $0.2$ and
$0.5$: about $0.3$ for metals, close to $0.5$ for rubber, and nearly $0$ for cork.
\end{defbox}

Why does $0.5$ appear as an upper value? Because of what stretching does to the volume.

\begin{derivbox}{Volume change of a stretched rod, and why $\nu\le\tfrac12$}
\stepn{1} For a rod of square cross-section $d\times d$, the volume is $V=\ell d^{2}$. Take logarithms and
differentiate (the logarithmic differentiation of \chref{ch:toolkit}):
\[
\frac{\Delta V}{V}=\frac{\Delta\ell}{\ell}+2\frac{\Delta d}{d}.
\]
(A circular section gives the same result, since its area is also proportional to $d^{2}$.)

\stepn{2} Substitute $\Delta d/d=-\nu\,\Delta\ell/\ell$:
\[
\boxed{\;\frac{\Delta V}{V}=(1-2\nu)\,\frac{\Delta\ell}{\ell}\;}
\]
\stepn{3} Read it. For $\nu=\tfrac12$ the volume does not change at all: the material flows sideways
exactly enough to make room. For $\nu<\tfrac12$, stretching increases the volume. A value above
$\tfrac12$ would mean that pulling a rod \emph{reduces} its volume; by the relation $Y=3B(1-2\nu)$ quoted
below, it would also make the bulk modulus negative, so that squeezing the material would make it expand,
and no stable material behaves like that. Rubber, with $\nu\approx0.5$, is the nearly volume-conserving
extreme, which is why its bulk modulus is so much larger than its Young's modulus.
\end{derivbox}

\inwords{Poisson's ratio ties the sideways squeeze to the lengthwise pull. Cork ($\nu\approx0$) barely
changes width when pushed into a bottle neck, which is exactly why it makes a good stopper; rubber
($\nu\approx0.5$) bulges strongly and keeps its volume. For a material that behaves the same in every
direction the four constants $Y$, $G$, $B$ and $\nu$ are not independent: any two fix the others, through
$Y=2G(1+\nu)=3B(1-2\nu)$. These relations are stated here without proof, because the derivation needs the
full three-dimensional theory of strain; the steel values in Table~\ref{tab:solids-moduli} satisfy them
with $\nu\approx0.29$.}

\section{The Stress--Strain Curve}\label{sec:solids-curve}

Hooke's law holds only for small strains. Pull a real material further and it passes through a sequence
of regimes, recorded in its \textbf{stress--strain curve}, measured by stretching a sample slowly in a
testing machine and plotting stress against strain (Figure~\ref{fig:solids-stressstrain}).

\begin{figure}[htbp]
\centering
\tikzfig{solids_stressstrain}
\figcaption{fig:solids-stressstrain}{Left: a ductile metal (strain axis not to scale; the straight part
really ends near a strain of $0.002$). Hooke's law holds up to A; beyond the yield point B the deformation
becomes permanent, so unloading from C retraces a line parallel to the elastic one and leaves a permanent
set. Past the ultimate stress D the sample narrows locally (necks) and breaks at E. Right: rubber has no
straight region, stretches to several times its length, and returns along a lower path, so a loading and
unloading cycle turns the area between the curves into heat.}
\end{figure}

\begin{defbox}{Landmarks on the stress--strain curve}
\begin{itemize}
\item \textbf{Proportional limit} (A): up to here stress is proportional to strain, and the moduli
      describe the material.
\item \textbf{Elastic limit, or yield point} (B): up to here the material returns to its original
      length when unloaded, even if the last part of the curve is not quite straight. The stress at B is
      the \textbf{yield strength}, the practical limit for design.
\item \textbf{Plastic region} (B to D): the deformation is now partly permanent. Unloaded from a point
      such as C, the material retraces a line parallel to the elastic one and keeps a \textbf{permanent
      set}.
\item \textbf{Ultimate tensile strength} (D): the greatest stress the material can bear. Beyond it the
      sample narrows at one place (necking), and the force needed to stretch it further falls.
\item \textbf{Fracture point} (E): where it breaks.
\end{itemize}
A \textbf{ductile} material, such as copper or mild steel, has a long plastic region and can be drawn into
wire; a \textbf{brittle} one, such as glass or cast iron, breaks at or near its elastic limit with almost
no warning.
\end{defbox}

\begin{remarkbox}{Elastomers and elastic hysteresis}
Rubber, and the elastic tissue of the aorta, are \textbf{elastomers}: they stretch to several times their
length, never obey Hooke's law, and still return to their original length. Their unloading curve lies
below the loading curve (Figure~\ref{fig:solids-stressstrain}, right), so a full cycle of stretching and
release does net work on the material; that work, equal to the area of the loop, becomes heat. This
\textbf{elastic hysteresis} is useful in shock absorbers and the soles of running shoes, and it is why a
car tyre warms up as it rolls.
\end{remarkbox}

\why{The curve explains why the structures around us do not collapse. Engineers keep the working stress
of a lift cable, a bridge girder or an aircraft spar several times below the yield strength; the ratio
is the \textbf{factor of safety}, chosen to cover flaws, fatigue and loads nobody foresaw. The plastic
region is useful knowledge too: it is what a blacksmith exploits to forge a horseshoe, what lets a car's
crumple zone absorb a crash, and what makes soft metals drawable into wires and beatable into sheets.
\chref{ch:thermal} (Section~\ref{sec:thermal-stress}) will show another source of stress: a steel rail clamped at both ends and warmed by
\SI{40}{K} develops a stress of about \SI{96}{MPa}, a large fraction of the yield strength of ordinary
steel, which is why railway lines and bridges need expansion gaps.}

\begin{exbox}{Choosing a crane rope with a factor of safety}\label{ex:solids-crane}
A crane must lift loads of up to \SI{1.0e4}{kg}. Its steel rope has a yield strength of
\SI{3.0e8}{Pa}. Find the least radius of rope that would just reach yield, and the radius needed for a
factor of safety of $10$.

\stepn{1} The weight is $mg=\SI{9.8e4}{N}$. Yield is reached when $mg/(\pi r^{2})=\num{3.0e8}$, so
$r=\sqrt{mg/(\pi\sigma_y)}=\sqrt{\num{9.8e4}/(\pi\times\num{3.0e8})}=\SI{1.0}{cm}$.

\stepn{2} A factor of safety of $10$ means a working stress of one tenth of the yield strength, so the
area must be ten times larger and the radius $\sqrt{10}$ times larger: $r=\SI{3.2}{cm}$.

\smallskip
Real crane ropes are made of many thin strands twisted together, which makes them flexible and means one
broken strand is not a disaster, but the area they must provide is set by exactly this calculation.
\ans{\SI{1.0}{cm}; \SI{3.2}{cm}}
\end{exbox}

\section{Elastic Potential Energy}\label{sec:solids-energy}

Stretching a wire takes work, and while the wire stays within its elastic limit that work is stored, to
be returned when the load is removed. \chref{ch:energy} found the energy of a stretched spring,
$U=\tfrac12kx^{2}$, by integrating a force that grows from zero. The same calculation for a wire gives a
result that belongs to the material.

\begin{derivbox}{Energy stored in a stretched wire}
\stepn{1} At extension $x$ the wire pulls back with force $F(x)=kx=\dfrac{YA}{\ell}\,x$. Stretching it a
further $dx$ takes work $F(x)\,dx$.

\stepn{2} Integrate from zero to the final extension $\Delta\ell$:
\[
U=\int_0^{\Delta\ell}\frac{YA}{\ell}\,x\,dx=\frac{YA}{\ell}\cdot\frac{(\Delta\ell)^{2}}{2}
=\frac12\,\underbrace{\frac{YA\,\Delta\ell}{\ell}}_{F}\,\Delta\ell=\frac12\,F\,\Delta\ell ,
\]
which is $\tfrac12kx^{2}$ with $k=YA/\ell$.

\stepn{3} Divide by the volume $A\ell$ to obtain the \textbf{energy density}, the energy stored per unit
volume, which no longer depends on the size of the wire:
\[
\boxed{\;u=\frac{U}{A\ell}=\frac12\,\sigma\,\epsilon=\frac12\,Y\epsilon^{2}=\frac{\sigma^{2}}{2Y}\;}
\]
\end{derivbox}

\inwords{The stored energy is $\tfrac12F\,\Delta\ell$, not $F\,\Delta\ell$, for the same reason as in a
spring's $\tfrac12kx^{2}$: the force grew linearly from zero, so its average over the stretch was half
its final value. Per unit volume the energy is half the stress times the strain, the area of the triangle
under the straight part of the stress--strain curve. Beyond the elastic limit the area under the curve is
still the work done, but only part of it comes back; the rest is the work of permanent deformation, which
ends up as heat.}

\begin{exbox}{Energy density in a loaded steel wire}\label{ex:solids-energy}
Find the elastic energy per unit volume in the steel wire of Example~\ref{ex:solids-wire}, and the total
energy it stores.

\stepn{1} $u=\dfrac{\sigma^{2}}{2Y}=\dfrac{(\num{9.8e7})^{2}}{2(\num{2.0e11})}=\SI{2.4e4}{J/m^3}$.

\stepn{2} The volume is $A\ell=\SI{1.0e-6}{m^3}$, so $U=\SI{0.024}{J}$. Check:
$\tfrac12F\Delta\ell=\tfrac12(98)(\num{4.9e-4})=\SI{0.024}{J}$.
\ans{\SI{2.4e4}{J/m^3}; \SI{24}{mJ}}
\end{exbox}

\section{Applications of Elastic Behaviour}\label{sec:solids-apps}

The point of understanding elasticity is to build with it: to choose shapes and materials that carry
loads without sagging too far, breaking, or wasting material. Three classic decisions follow from the
ideas above, and each comes with an honest statement of what the model leaves out.

\subsection{Beams: why girders are I-shaped and bones are hollow}

A horizontal beam bends under a load, its upper layers stretching and its lower layers compressing
(Figure~\ref{fig:solids-beam}). How far it sags depends on the material through $Y$ and on the shape of
its cross-section through a single geometric number.

\begin{figure}[htbp]
\centering
\tikzfig{solids_beam}
\figcaption{fig:solids-beam}{Left: a cantilever (sag exaggerated). Above the neutral axis the layers
stretch, below it they compress, and the axis itself keeps its length. Right: four cross-sections of
\emph{equal area}, so equal mass per metre, with their resistance to bending (the second moment of area
$I_g$) relative to the square. Moving material away from the neutral axis, by standing a plank on edge,
gathering it into flanges or rolling it into a tube, makes the same material several times stiffer.}
\end{figure}

\begin{defbox}{The depression of a cantilever}
A horizontal beam of length $\ell$, clamped at one end and loaded with a force $F$ at the free end, sags at
the tip by a \textbf{depression}
\[
\delta=\frac{F\ell^{3}}{3\,Y\,I_g},\qquad I_g=\int y^{2}\,dA ,
\]
where $Y$ is Young's modulus and $I_g$ is the \textbf{second moment of area} of the cross-section about its
horizontal neutral axis ($y$ is the height of an element of area $dA$ above that axis). The integral has
the same form as the moment of inertia $I=\int r^{2}\,dm$ of \chref{ch:rotation}, with area in place of
mass, which is why it rewards placing material \emph{far} from the axis. For a rectangle of width $b$ and
depth $d$, $I_g=bd^{3}/12$. The result is derived just below; three consequences matter more than the
derivation:
\begin{itemize}
\item \textbf{The $\ell^{3}$ is severe.} Doubling the length of a beam makes it sag eight times as much.
\item \textbf{Depth beats width.} Since $I_g\propto bd^{3}$, a plank laid on edge is far stiffer than the
      same plank laid flat: standing it up, so that depth and width exchange, multiplies the stiffness by
      $(d/b)^{2}$. That is why floor joists stand on edge.
\item \textbf{Shape beats bulk.} An I-section (material gathered into top and bottom flanges) and a hollow
      tube are several times stiffer than a solid bar of the same mass (Figure~\ref{fig:solids-beam}),
      which is why steel girders and rails are I-shaped, and why bones, bamboo and bicycle frames are
      hollow.
\end{itemize}
\end{defbox}

\begin{derivbox}{Where $\delta=F\ell^{3}/3YI_g$ comes from (beyond the syllabus)}
This derivation is above the level of the examination, but it uses only Hooke's law, the turning effect
of a force and two integrations, and it shows exactly why shape enters through $\int y^{2}\,dA$. Read it
for understanding, not for reproduction.

\stepn{1} \textbf{Bending stretches one side and compresses the other.} Bend a short piece of beam into an
arc of radius $R$. The neutral axis keeps its length; a layer a height $y$ above it lies on an arc of
radius $R+y$, so it is stretched by the strain $y/R$, proportional to its distance from the axis. By
Hooke's law the stress there is $\sigma=Yy/R$.

\stepn{2} \textbf{Add up the layers into a bending moment.} The force on a strip of area $dA$ is
$\sigma\,dA$, at lever arm $y$ about the neutral axis, so its turning effect (its moment, as for the torque
of \chref{ch:rotation}) is $y\,\sigma\,dA$. Summing over the cross-section,
\[
M=\int y\,\sigma\,dA=\frac{Y}{R}\int y^{2}\,dA=\frac{YI_g}{R}.
\]
Because $y$ is \emph{squared}, material far from the axis counts for far more. Rearranged, this is the
\textbf{bending equation}, $1/R=M/(YI_g)$.

\stepn{3} \textbf{The moment along a cantilever.} At a distance $x$ from the wall, the part of the beam
beyond $x$ carries the load $F$ at a further distance $(\ell-x)$, so the moment there is $M(x)=F(\ell-x)$.
For a small sag $y_s(x)$ (measured downwards) the curvature is $1/R=d^{2}y_s/dx^{2}$, a result about
curves quoted here from mathematics, so
\[
YI_g\,\frac{d^{2}y_s}{dx^{2}}=F(\ell-x).
\]
\stepn{4} \textbf{Integrate twice.} At the clamped end the sag and the slope are both zero, which fixes
both constants of integration at zero:
\[
YI_g\,\frac{dy_s}{dx}=F\Big(\ell x-\frac{x^{2}}{2}\Big),\qquad
YI_g\,y_s=F\Big(\frac{\ell x^{2}}{2}-\frac{x^{3}}{6}\Big).
\]
At the free end, $x=\ell$, the bracket is $\tfrac12\ell^{3}-\tfrac16\ell^{3}=\tfrac13\ell^{3}$, so
\[
\boxed{\;\delta=y_s(\ell)=\frac{F\ell^{3}}{3YI_g}\;}
\]
\stepn{5} \textbf{A beam supported at both ends.} A beam of length $\ell$ resting on two supports and
loaded by $F$ at its centre is, by symmetry, level at the centre, so each half behaves like a cantilever of
length $\ell/2$ clamped at the middle and pushed up by $F/2$ at its support. Its sag is therefore
\[
\delta=\frac{(F/2)(\ell/2)^{3}}{3YI_g}=\frac{F\ell^{3}}{48\,YI_g}=\frac{F\ell^{3}}{4\,Y\,bd^{3}}
\quad\text{for a rectangular section,}
\]
the form in which the result usually appears in examinations.
\end{derivbox}

\begin{exbox}{A tube against a solid rod}\label{ex:solids-tube}
A solid circular rod of radius $r$ and a tube of outer radius $2r$ are made of the same material and have
the same mass per metre. Compare their resistance to bending. ($I_g=\tfrac14\pi R^{4}$ for a solid disc of
radius $R$ about a diameter, stated here from the integral $\int y^2\,dA$.)

\stepn{1} Equal areas: $\pi\big((2r)^{2}-r_{\text{in}}^{2}\big)=\pi r^{2}$, so the inner radius is
$r_{\text{in}}=\sqrt3\,r$.

\stepn{2} The tube's $I_g$ is that of the full disc minus that of the hole:
$\tfrac14\pi\big((2r)^{4}-(\sqrt3\,r)^{4}\big)=\tfrac14\pi(16-9)r^{4}=7\times\tfrac14\pi r^{4}$.

\smallskip
The tube is seven times stiffer than the solid rod for the same material and weight. It is weaker in one
respect the model ignores: a very thin-walled tube can buckle or dent locally, which is why bicycle tubes
and bones are not made arbitrarily thin. \ans{tube $7\times$ stiffer}
\end{exbox}

\subsection{Self-weight: how high can a mountain be?}

A structure must carry its own weight, and the stress at its base grows with its height. That sets an
upper limit on the height of anything built of a given material, including mountains.

\begin{exbox}{The tallest possible mountain}\label{ex:solids-mountain}
Model a mountain as a column of rock of uniform cross-section and height $h$. The stress at its base is
the weight above divided by the area, $\rho Ahg/A=\rho gh$. Rock begins to flow plastically at a stress of
about \SI{3e8}{Pa}; its density is about \SI{3e3}{kg/m^3}. Estimate the greatest height.

\stepn{1} Set $\rho gh=\sigma_y$: $h=\dfrac{\sigma_y}{\rho g}=\dfrac{\num{3e8}}{(\num{3e3})(9.8)}\approx\SI{10}{km}$.

\smallskip
The estimate is crude (a real mountain is wider at the base, which lowers the stress, and rock strength
varies), but it is the right size: Mount Everest is \SI{8.8}{km} high, and no mountain on Earth is much
taller. On Mars, where $g$ is about $0.38$ of its value on Earth, Olympus Mons rises over \SI{20}{km},
as the same estimate predicts. \ans{about \SI{10}{km}}
\end{exbox}

\why{Elasticity is where physics meets civil and mechanical engineering. The thickness of a crane rope is
set so that its working stress is a safe fraction of its yield strength; the ropes, pillars and beams of
every structure are sized by $\Delta\ell=F\ell/AY$, the stress limit, and the depression formulas of this
chapter; and the choice to alloy iron with a little carbon, or to prefer aluminium for aircraft, trades
stiffness, strength and weight. The next chapter turns from solids, which resist a change of shape, to
fluids, which do not.}

\section{Summary}\label{sec:summary-solids}

Stress (force per area) and strain (fractional deformation) describe a material independently of the size
of the sample; within the elastic limit they are proportional, and the constant of proportionality, one
for each kind of deformation, is the modulus. A wire is a spring whose constant is the modulus times
$A/\ell$, and the energy it stores is half the stress times the strain per unit volume. Beyond the elastic
limit a material deforms permanently and eventually breaks, and the design of every structure keeps its
stresses well inside that limit, using shape as much as material to do so.

\begin{itemize}
\item Stress $\sigma=F/A$ (\si{Pa}); longitudinal strain $\epsilon=\Delta\ell/\ell$; shear strain
      $\theta\approx\Delta x/h$; volume strain $\Delta V/V$.
\item Moduli: $Y=\sigma/\epsilon$, $G=(F/A)/\theta$, $B=-\Delta P/(\Delta V/V)$; fluids have only $B$.
\item $\Delta\ell=F\ell/AY$, and a wire is a spring with $k=YA/\ell$; the stretch is proportional to
      $\ell/r^{2}$.
\item Poisson's ratio $\nu=-(\Delta d/d)/(\Delta\ell/\ell)$; a stretched rod changes volume by
      $\Delta V/V=(1-2\nu)\,\Delta\ell/\ell$, so $\nu\le\tfrac12$.
\item Stress--strain curve: proportional limit, elastic limit (yield), plastic region with permanent set,
      ultimate strength, fracture. Breaking depends on stress, never on length.
\item Elastic energy $U=\tfrac12F\,\Delta\ell$; energy density $u=\tfrac12\sigma\epsilon=\sigma^{2}/2Y$.
\item A cantilever sags by $F\ell^{3}/3YI_g$ with $I_g=\int y^{2}dA$; depth and distance from the neutral
      axis buy stiffness, which is why girders are I-shaped and bones hollow.
\end{itemize}

\section{Problem Types}\label{sec:types-solids}

Six types cover the elasticity questions examiners set. Almost all of them rest on $\Delta\ell=F\ell/AY$
and on remembering that a material fails at a given \emph{stress}.

\archq{Elongation when the dimensions change}
{``A wire's length, radius or load is changed; find the new extension,'' or ``compare two wires of the
same material.''}
{Use $\Delta\ell=\dfrac{F\ell}{AY}=\dfrac{F\ell}{\pi r^{2}Y}$ and track each factor: $\Delta\ell\propto F$,
$\propto\ell$, $\propto1/r^{2}$. If the volume is fixed, a change of length forces a change of area.}
{The radius enters \emph{squared}: halving the radius quadruples the stretch. And ``same volume, twice the
length'' halves the area as well.}
{A wire stretches by \SI{1.0}{mm} under a certain load. A second wire of the same material and the same
volume, but twice the length, carries the same load. Its extension is:}
{\SI{1.0}{mm}}{\SI{2.0}{mm}}{\SI{4.0}{mm}}{\SI{8.0}{mm}\archans{C}{Same volume with twice the length
means half the area. Since $\Delta\ell\propto\ell/A$, the extension is multiplied by $2\times2=4$, giving
\SI{4.0}{mm}. Option (B) forgets the change of area; option (D) treats the area as falling by a factor of
four, as if the radius (not the area) had halved.}}

\archq{The compound wire}
{Two wires of different materials joined end to end (in series) or side by side (in parallel); the total
extension or the share of the load is wanted.}
{In series the same force acts through both and the extensions add. In parallel the extension is common
and the load divides in proportion to each wire's $YA/\ell$, its spring constant.}
{In series the \emph{force} is common, not the stress (the areas may differ); in parallel the
\emph{extension} is common. Decide first what is shared.}
{Two wires of equal length and equal area, with Young's moduli $Y$ and $2Y$, are joined end to end and
pulled. The ratio of their extensions (first to second) is:}
{$1:1$}{$2:1$}{$1:2$}{$4:1$\archans{B}{In series the same force acts on both, and with equal areas the
stresses are equal too. Since $\Delta\ell\propto1/Y$, the stiffer wire stretches half as much: $2:1$.
Option (C) inverts the dependence on $Y$; option (A) assumes equal extensions, which is the parallel
case.}}

\archq{Elastic energy stored}
{``A wire is stretched by $\Delta\ell$ (or to a stress $\sigma$); find the energy stored,'' or ``compare
energies when the dimensions change.''}
{$U=\tfrac12F\,\Delta\ell=\tfrac12\dfrac{YA}{\ell}(\Delta\ell)^{2}$; per unit volume
$u=\tfrac12\sigma\epsilon=\tfrac12Y\epsilon^{2}$.}
{The energy is $\tfrac12F\,\Delta\ell$, not $F\,\Delta\ell$, because the force grew from zero; and it goes
as the \emph{square} of the extension.}
{A wire stretched by \SI{1}{mm} stores energy $U$. Stretched instead by \SI{2}{mm}, still within the
elastic limit, it stores:}
{$U$}{$2U$}{$U/2$}{$4U$\archans{D}{$U=\tfrac12\dfrac{YA}{\ell}(\Delta\ell)^{2}\propto(\Delta\ell)^{2}$, so
doubling the extension quadruples the energy. Option (B) treats the energy as proportional to the
extension, forgetting that the force doubles too.}}

\archq{Breaking load and factor of safety}
{``Will a wire break? Find the greatest load, or the factor of safety,'' given a breaking stress or yield
strength.}
{Failure depends on stress: $F_{\max}=\sigma_{\text{break}}A$, with $A=\pi r^{2}=\pi d^{2}/4$. The factor of
safety is $\sigma_{\text{break}}/\sigma_{\text{working}}$.}
{Using the diameter as the radius, forgetting to divide the force by $g$ to get a mass, or thinking the
length matters.}
{A steel wire of \emph{diameter} \SI{1.0}{mm} has a breaking stress of \SI{4.0e8}{Pa}. The greatest mass
it can support is about ($g=\SI{9.8}{m/s^2}$):}
{\SI{32}{kg}}{\SI{128}{kg}}{\SI{41}{kg}}{\SI{314}{kg}\archans{A}{$A=\pi(\SI{0.5e-3}{m})^{2}=\SI{7.85e-7}{m^2}$,
so $F_{\max}=\sigma A=\SI{314}{N}$ and $m=F_{\max}/g=\SI{32}{kg}$. Option (B) uses the diameter as the
radius; option (C) takes the area as $d^{2}$; option (D) forgets to divide by $g$.}}

\archq{Bulk modulus and compression}
{A body is placed under extra pressure on all sides (deep water, a hydraulic press); a change in volume,
radius or density is wanted.}
{$\Delta V/V=-\Delta P/B$. A linear dimension changes by one third of the volume strain, since
$V\propto r^{3}$; the density rises by the same fraction as the volume falls.}
{Reporting the volume strain when the change in radius is asked for, or multiplying by $3$ where one
should divide.}
{A steel sphere ($B=\SI{1.6e11}{Pa}$) is lowered into the deep ocean, where the pressure exceeds that at
the surface by \SI{1.0e8}{Pa}. Its radius decreases by about:}
{$0.063\%$}{$0.19\%$}{$0.021\%$}{$0.031\%$\archans{C}{$\Delta V/V=\num{1.0e8}/\num{1.6e11}=\num{6.25e-4}$,
and $\Delta r/r=\tfrac13\,\Delta V/V=\num{2.1e-4}=0.021\%$. Option (A) is the volume strain itself; option
(B) multiplies by $3$ instead of dividing; option (D) divides by $2$.}}

\archq{Poisson's ratio and volume change}
{A wire or rod is stretched and the change in its diameter or volume is wanted, with $\nu$ given.}
{$\Delta d/d=-\nu\,\Delta\ell/\ell$ and $\Delta V/V=(1-2\nu)\,\Delta\ell/\ell$.}
{Assuming the volume stays constant (true only for $\nu=\tfrac12$), or forgetting that the lateral strain
enters twice because there are two sideways directions.}
{A wire with Poisson's ratio $0.25$ is stretched to a longitudinal strain of $0.20\%$. Its volume:}
{increases by $0.10\%$}{increases by $0.20\%$}{does not change}{increases by $0.30\%$\archans{A}{$\Delta
V/V=(1-2\times0.25)(0.20\%)=0.10\%$. Option (B) ignores the sideways contraction; option (C) assumes
volume is conserved; option (D) adds the lateral strain instead of subtracting it.}}

\section{Practice Problems}\label{sec:practice-solids}

The problems below are of standard examination difficulty. Take $g=\SI{9.8}{m/s^2}$ and use
Table~\ref{tab:solids-moduli} where a modulus is not given.

\begin{enumerate}[itemsep=3pt]
\item A copper wire \SI{2.0}{m} long and \SI{1.0}{mm} in diameter carries a \SI{2.0}{kg} load. Find its
      extension.
\item Two wires of the same material have lengths in the ratio $1:2$ and diameters in the ratio $2:1$.
      They carry equal loads. Find the ratio of their extensions.
\item What extra pressure compresses water by $0.10\%$ of its volume? Express it also in atmospheres
      ($\SI{1}{atm}=\SI{1.013e5}{Pa}$, as in \chref{ch:fluids}).
\item A steel wire is stretched to a strain of $0.10\%$. Find the elastic energy stored per cubic metre.
\item A rubber cord with Poisson's ratio $0.50$ is stretched by $2.0\%$. By what percentage does its
      diameter change, and what happens to its volume?
\item A steel wire and a brass wire ($Y=\SI{1.0e11}{Pa}$), each \SI{1.0}{m} long with cross-section
      \SI{1.0}{mm^2}, are joined end to end and hung from a ceiling with a \SI{10}{kg} load at the bottom.
      Neglecting the weight of the wires, find the total extension.
\item A steel wire hangs vertically from a support. Its breaking stress is \SI{4.0e8}{Pa} and its density
      \SI{7800}{kg/m^3}. How long can it be before it breaks under its own weight? Does the answer depend
      on its thickness?
\end{enumerate}

\subsection*{Answers}
{\small
\begin{enumerate}[itemsep=2pt]
\item $A=\pi(\num{0.5e-3})^{2}=\SI{7.85e-7}{m^2}$; $\Delta\ell=\dfrac{(19.6)(2.0)}{(\num{7.85e-7})(\num{1.2e11})}
      =\SI{4.2e-4}{m}\approx\SI{0.42}{mm}$.
\item $\Delta\ell\propto\ell/d^{2}$: $\dfrac{1/2^{2}}{2/1^{2}}=\dfrac18$, so $1:8$.
\item $\Delta P=B\,|\Delta V/V|=(\num{2.2e9})(\num{1.0e-3})=\SI{2.2e6}{Pa}$, about $22$ atmospheres.
\item $u=\tfrac12Y\epsilon^{2}=\tfrac12(\num{2.0e11})(\num{1.0e-3})^{2}=\SI{1.0e5}{J/m^3}$.
\item $\Delta d/d=-\nu\,\Delta\ell/\ell=-1.0\%$: the diameter shrinks by $1.0\%$. With $\nu=\tfrac12$,
      $\Delta V/V=(1-2\nu)\,\Delta\ell/\ell=0$, so the volume is unchanged (to first order in the strain).
\item The same force, $F=\SI{98}{N}$, acts through both wires, so the extensions add:
      $\dfrac{F\ell}{A}\Big(\dfrac1{Y_s}+\dfrac1{Y_b}\Big)=\num{9.8e7}\,(\num{5e-12}+\num{1e-11})
      =\SI{1.47e-3}{m}\approx\SI{1.5}{mm}$, of which the brass provides two thirds.
\item The top of the wire carries the weight of the whole wire, $\rho A\ell g$, so the greatest stress is
      $\rho g\ell$. Setting it equal to the breaking stress, $\ell=\dfrac{\num{4.0e8}}{(7800)(9.8)}
      \approx\SI{5.2}{km}$. The area cancels, so the thickness does not matter: a thicker wire is heavier in
      exact proportion to its greater strength. (This is the mountain estimate of
      Example~\ref{ex:solids-mountain} for a wire.)
\end{enumerate}
}

\section{Challenge Problems}\label{sec:challenge-solids}

The five problems below go beyond examination level, chiefly because each needs a stress or strain that
varies from place to place, or a quantity from another chapter. Attempt each fully before reading its
solution.

\begin{challengebox}{A rod spun about its end (builds on stress and strain)}
A uniform rod of length $\ell=\SI{1.0}{m}$, density $\rho=\SI{7800}{kg/m^3}$ and Young's modulus
$Y=\SI{2.0e11}{Pa}$ is spun in a horizontal plane about a vertical axis through one end at
$\omega=\SI{100}{rad/s}$ (Figure~\ref{fig:solids-spinrod}). Find the greatest tensile stress in the rod
and its total extension. (Neglect gravity.)
\challengeans{The tension is not uniform. The part of the rod beyond a distance $x$ from the axis moves
in circles, and the only horizontal force on it is the tension $T(x)$ at the cut, which must provide its
total centripetal force (\chref{ch:circdyn}). A slice of length $dr$ at radius $r$ has mass $\rho A\,dr$
and needs the force $\rho A\,dr\,\omega^{2}r$, so
\[
T(x)=\int_x^{\ell}\rho A\omega^{2}r\,dr=\tfrac12\rho A\omega^{2}(\ell^{2}-x^{2}),
\]
largest at the axis. The greatest stress is
\[
\sigma_{\max}=\frac{T(0)}{A}=\tfrac12\rho\omega^{2}\ell^{2}=\tfrac12(7800)(100)^{2}(1.0)^{2}
\approx\SI{39}{MPa}.
\]
The extension adds up the local strains $\sigma(x)/Y$ over the length:
\[
\Delta\ell=\int_0^{\ell}\frac{\sigma(x)}{Y}\,dx=\frac{\rho\omega^{2}}{2Y}\int_0^{\ell}(\ell^{2}-x^{2})\,dx
=\frac{\rho\omega^{2}\ell^{3}}{3Y}\approx\SI{0.13}{mm}.
\]
The $\tfrac13$, where a uniform tension would give $\tfrac12$, is the signature of a tension that varies
quadratically along the rod.
\ans{\SI{39}{MPa}; \SI{0.13}{mm}}}
\end{challengebox}

\begin{figure}[htbp]
\centering
\tikzfig{solids_spinrod}
\figcaption{fig:solids-spinrod}{A rod spun about one end, seen from above. The cut at distance $x$ must
pull the whole outer part towards the axis, so the tension $T(x)$ is the total centripetal force needed
beyond $x$: greatest at the axis, which carries the whole rod, and zero at the free end, which carries
nothing. The graph shows $T\propto\ell^{2}-x^{2}$.}
\end{figure}

\begin{challengebox}{A wire stretched by its own weight (builds on stress and strain)}
A uniform steel wire of length $\ell=\SI{100}{m}$ hangs vertically from one end ($\rho=\SI{7800}{kg/m^3}$,
$Y=\SI{2.0e11}{Pa}$). Find its extension under its own weight. Why is it not $\rho g\ell^{2}/Y$?
\challengeans{The tension is not uniform. At a height $x$ above the free end the wire supports only the
weight of the length $x$ below it, $\rho Axg$, so the stress there is $\sigma(x)=\rho gx$: largest at the
top, zero at the bottom. An element $dx$ stretches by $(\sigma/Y)\,dx$, so
\[
\Delta\ell=\int_0^{\ell}\frac{\rho gx}{Y}\,dx=\frac{\rho g\ell^{2}}{2Y}
=\frac{(7800)(9.8)(100)^{2}}{2(\num{2.0e11})}\approx\SI{1.9}{mm}.
\]
The factor $\tfrac12$ comes from averaging a stress that rises linearly from bottom to top: the wire
stretches as if its whole weight hung from its midpoint. The same averaging gives the force on a dam wall
in \chref{ch:fluids}, where the pressure rises linearly with depth.
\ans{\SI{1.9}{mm}}}
\end{challengebox}

\begin{challengebox}{Twisting a wire (builds on the shear modulus)}
A wire of radius $r=\SI{0.50}{mm}$, length $\ell=\SI{1.0}{m}$ and shear modulus $G=\SI{8.0e10}{Pa}$ is
clamped at its upper end. What torque must be applied to the lower end to twist it through a right angle?
\challengeans{Divide the wire into thin cylindrical shells of radius $r'$ and thickness $dr'$. When the free
end turns through the angle $\psi$, a line drawn along the surface of a shell is tilted through the shear
angle $r'\psi/\ell$, so the shear stress in the shell is $Gr'\psi/\ell$. That stress acts over the shell's
end area $2\pi r'\,dr'$ at lever arm $r'$, giving a torque $d\tau=(Gr'\psi/\ell)(2\pi r'\,dr')\,r'$. Adding the
shells,
\[
\tau=\frac{2\pi G\psi}{\ell}\int_0^{r}r'^{3}\,dr'=\frac{\pi Gr^{4}}{2\ell}\,\psi
=\frac{\pi(\num{8.0e10})(\num{0.50e-3})^{4}}{2(1.0)}\cdot\frac{\pi}{2}\approx\SI{12}{mN.m}.
\]
Note the $r^{4}$: halving the radius makes the wire sixteen times easier to twist. That is why a very thin
fibre is used as the suspension of a torsion balance, the instrument with which the gravitational constant
of \chref{ch:gravitation} was first measured.
\ans{\SI{12}{mN.m}}}
\end{challengebox}

\begin{challengebox}{Two rods clamped end to end (builds on Young's modulus)}
A steel rod and an aluminium rod, each \SI{0.50}{m} long and of the same cross-section, are joined end to
end and fixed between two immovable walls at \SI{20}{\celsius}, then heated to \SI{70}{\celsius}. Free to
expand, a rod would lengthen by the fraction $\alpha_\ell\,\Delta T$, where $\alpha_\ell$ is its linear
expansion coefficient (defined in \chref{ch:thermal}) and $\Delta T$ the rise in temperature. Take
$\alpha_\ell=\SI{1.2e-5}{K^{-1}}$ and $Y=\SI{2.0e11}{Pa}$ for steel, $\alpha_\ell=\SI{2.4e-5}{K^{-1}}$
and $Y=\SI{7.0e10}{Pa}$ for aluminium. Find the stress in the rods, and how far, and which way, the
junction moves.
\challengeans{A single rod clamped between walls is the simplest case: the walls impose a compressive
strain equal to the free expansion, so $\sigma=Y\alpha_\ell\,\Delta T$, independent of length and area
(derived in full in \chref{ch:thermal}, Section~\ref{sec:thermal-stress}). Two rods in series need one
more idea. The force is the same throughout (each rod pushes on the other), and with equal areas so is
the stress $\sigma$. The walls fix the \emph{total} length, so the two free expansions must be undone by
the two elastic compressions together:
\[
\ell(\alpha_{\text{s}}+\alpha_{\text{a}})\Delta T=\sigma\ell\Big(\frac{1}{Y_{\text{s}}}+\frac{1}{Y_{\text{a}}}\Big),
\quad\text{so}\quad
\sigma=\frac{(\alpha_{\text{s}}+\alpha_{\text{a}})\Delta T}{1/Y_{\text{s}}+1/Y_{\text{a}}}
=\frac{(\num{3.6e-5})(50)}{\num{1.93e-11}}\approx\SI{93}{MPa}
\]
(compressive), with subscripts s for steel and a for aluminium.
The steel's net strain is $\alpha_{\text{s}}\Delta T-\sigma/Y_{\text{s}}=\num{6.0e-4}-\num{4.7e-4}
=+\num{1.3e-4}$: it ends up slightly \emph{longer}. The aluminium's is
$\num{1.2e-3}-\num{1.33e-3}=-\num{1.3e-4}$: it is squeezed shorter than it started, by exactly the amount
the steel grew. So the junction moves towards the aluminium, by $(0.50)(\num{1.3e-4})\approx\SI{0.07}{mm}$.
The softer, more expansive aluminium loses the contest. The stress lies between the values the two rods
would reach clamped alone (\SI{120}{MPa} for steel, \SI{84}{MPa} for aluminium), and it does not depend on
the common length, which is why rails and bridges need expansion gaps however long or short they are.
\ans{\SI{93}{MPa}; the junction moves \SI{0.07}{mm} towards the aluminium}}
\end{challengebox}

\begin{challengebox}{Where does half the energy go? (builds on elastic energy)}
A wire of natural length \SI{1.0}{m} and cross-section \SI{1.0}{mm^2} ($Y=\SI{2.0e11}{Pa}$) supports a
\SI{10}{kg} load. Find the elastic energy stored, and compare it with the work done by gravity as the load
descends through the extension.
\challengeans{The extension is $\Delta\ell=\dfrac{mg\ell}{AY}=\SI{0.49}{mm}$ (Example~\ref{ex:solids-wire}),
and the wire stores
\[
U=\tfrac12F\,\Delta\ell=\tfrac12(98)(\num{0.49e-3})\approx\SI{24}{mJ}.
\]
But gravity does $W=mg\,\Delta\ell=\SI{48}{mJ}$ as the load descends, twice as much. Where does the other
half go? It depends on how the load is attached. If it is released suddenly from the unstretched position,
it overshoots, oscillates about the equilibrium and is brought to rest by internal friction, which turns
the surplus \SI{24}{mJ} into heat. If instead it is lowered gently by a hand, the hand pushes up throughout
and does work $-\SI{24}{mJ}$, so the net work on the load is just what the wire stores. Either way the wire
holds exactly half the work done by gravity: the energy method of \chref{ch:energy} must account for every
force that does work.
\ans{\SI{24}{mJ}, half of the \SI{48}{mJ} done by gravity}}
\end{challengebox}
